\documentclass[12pt]{article}
\usepackage[margin=1in]{geometry}
\usepackage[UTF8,fontset=fandol]{ctex}
\renewcommand{\figurename}{Figure}
\renewcommand{\tablename}{Table}
\usepackage{amsmath, amssymb, graphicx, setspace, natbib}
\usepackage{booktabs,tabularx}
\usepackage{multirow}
\usepackage{comment}
\usepackage{enumitem}
\usepackage{array}

\title{Tidal Lanes: Shenzhen Data Pipeline Draft}
\author{Xiaoruo Feng}
\date{April 2026}

\begin{document}
\maketitle

\section{Data}
\subsection{Commuting-Flow Data}

The Shenzhen commuting matrix used in this draft is the 2024 100-meter grid origin--destination dataset. After standardizing field names for the stage-based pipeline, the sample contains 3,970,569 origin--destination pairs, 94,670 unique origin grids, 105,322 unique destination grids, and 110,339 unique grids in the union of origins and destinations.

The pipeline-level commuting file keeps grid identifiers, grid-center coordinates, and commuter counts. The transport-mode decomposition fields are not used in the current Shenzhen run. This is a deliberate design choice for this draft version. The downstream network and QSM stages are built from total commuters.

\subsection{Road Segment Speed Data}

Road speeds are built from the Raw\_fromBaidu source for Shenzhen. We match the Shenzhen road segment list to the high-frequency speed index by \texttt{roadseg\_id}, then retain the common support. The resulting speed table contains 10,435,104 records and 10,433 unique road segments.

The retained raw speed fields are \texttt{cityname}, \texttt{roadseg\_id}, \texttt{roadname}, \texttt{semantic}, \texttt{location}, \texttt{geometry}, \texttt{speed}, \texttt{day}, and \texttt{hour}. The temporal resolution remains ten-minute bins.

\section{Pipeline Overview}
\label{sec:sz_pipeline}

The Shenzhen pipeline run in this draft is \texttt{shenzhen\_test\_line\_v1\_run1}. The stage sequence is
\[
\texttt{stage01} \rightarrow \texttt{stage02} \rightarrow \texttt{stage03} \rightarrow \texttt{stage03b} \rightarrow \texttt{stage04} \rightarrow \texttt{stage05} \rightarrow \texttt{stage06} \rightarrow \texttt{stage07} \rightarrow \texttt{stage08} \rightarrow \texttt{stage09} \rightarrow \texttt{stage10}.
\]

Stage01 and stage02 construct the directed centerline backbone and the raw-to-centerline crosswalk. Stage03 and stage03b aggregate directional speeds and compute asymmetry diagnostics. Stage05 to stage08 build the square, hex, and Voronoi grid systems, route-level travel costs, OD reachability objects, and QSM-ready node-edge-OD inputs. Stage10 generates diagnostic figures.

\begin{figure}[h!]
\centering
\includegraphics[width=0.92\textwidth]{../../data_work/outputs/shenzhen_test_line_v1_run1/figures/map_raw_roads_15km.png}
\caption{Raw Shenzhen roads in the local diagnostic window used by the pipeline.}
\label{fig:sz_raw_roads_15km}
\end{figure}

\section{Road Network Construction and Centerline Matching}

\subsection{Stage01 centerline construction}

The Shenzhen road geometry enters stage01 through \texttt{raw\_data/gis/roads\_baidu/shenzhen\_roads.shp}, with the Shenzhen boundary file used in grid construction later. The skeleton branch runs with the retained production parameters in the project codebase, including 50-meter buffering and 5-meter raster resolution.

In the Shenzhen run, stage01 starts from 10,536 raw segments and outputs 12,714 undirected centerlines and 25,428 directed centerline links. After baseline filtering, 6,456 directed links remain in the retained backbone.

\begin{table}[h!]
\centering
\caption{Stage01 summary, Shenzhen run}
\label{tab:sz_stage01_summary}
\begin{tabular}{lr}
\toprule
Statistic & Value \\
\midrule
Raw segments & 10,536 \\
Undirected centerlines (all) & 12,714 \\
Directed centerlines (all) & 25,428 \\
Directed centerlines (baseline) & 6,456 \\
Total centerline length (km, all) & 2,420.4 \\
Baseline centerline length (km) & 2,356.0 \\
\bottomrule
\end{tabular}
\end{table}

\subsection{Stage02 matching quality}

Stage02 splits raw roads only where network geometry indicates a plausible multi-centerline overlap, then matches split segments to the directed backbone under the retained distance-cap rule. The Shenzhen run has 10,536 split segments, 9,909 baseline-kept split segments, and a baseline split match rate of 98.37\%. At the raw-edge level, 9,747 edges are matched and 162 are unmatched in the baseline-kept sample.

\begin{table}[h!]
\centering
\caption{Stage02 matching summary, Shenzhen run}
\label{tab:sz_stage02_summary}
\begin{tabular}{lr}
\toprule
Statistic & Value \\
\midrule
Split segments (all) & 10,536 \\
Split segments (baseline-kept) & 9,909 \\
Baseline split match rate & 98.37\% \\
Raw-edge match rate (baseline-kept) & 98.37\% \\
Matched split segments (baseline-kept) & 9,747 \\
Unmatched split segments (baseline-kept) & 162 \\
\bottomrule
\end{tabular}
\end{table}

\begin{figure}[h!]
\centering
\includegraphics[width=0.88\textwidth]{../../data_work/outputs/shenzhen_test_line_v1_run1/figures/map_raw_split_matched_vs_unmatched.png}
\caption{Raw split segments classified by matched and unmatched status in the Shenzhen run.}
\label{fig:sz_match_status_map}
\end{figure}

\begin{figure}[h!]
\centering
\includegraphics[width=0.82\textwidth]{../../data_work/outputs/shenzhen_test_line_v1_run1/figures/fig_length_cdf_4networks.png}
\caption{Length CDF comparison across retained Shenzhen network objects.}
\label{fig:sz_length_cdf}
\end{figure}

\section{Aggregating Raw Speeds to Directed Centerlines}

Stage03 maps Shenzhen speed observations to matched centerline segments and aggregates by weekday label and hour. The run yields 9,673,722 matched speed observations, 232,633 centerline-time rows, and 4,875 unique directed centerlines with valid speed support.

\begin{table}[h!]
\centering
\caption{Stage03 speed aggregation summary, Shenzhen run}
\label{tab:sz_stage03_summary}
\begin{tabular}{lr}
\toprule
Statistic & Value \\
\midrule
Matched speed observations & 9,673,722 \\
Centerline-time rows & 232,633 \\
Unique directed centerlines & 4,875 \\
Mean centerline speed (km/h) & 32.88 \\
Median centerline speed (km/h) & 29.64 \\
AM peak rows & 19,412 \\
PM peak rows & 19,453 \\
\bottomrule
\end{tabular}
\end{table}

\begin{figure}[h!]
\centering
\includegraphics[width=0.86\textwidth]{../../data_work/outputs/shenzhen_test_line_v1_run1/figures/hist_speed_raw_vs_centerline_all.png}
\caption{Speed distribution comparison between raw segment observations and aggregated centerline speeds.}
\label{fig:sz_speed_hist_raw_vs_centerline}
\end{figure}

\begin{figure}[h!]
\centering
\includegraphics[width=0.86\textwidth]{../../data_work/outputs/shenzhen_test_line_v1_run1/figures/diurnal_raw_centerline_combined.png}
\caption{Diurnal profile of raw and centerline speeds in the Shenzhen run.}
\label{fig:sz_speed_diurnal}
\end{figure}

\section{Directional Asymmetry and Tidal-Lane Screen}

Stage03b runs on the Shenzhen speed table produced by stage03. In the current run, the asymmetry output table is empty after the applied filters, so the baseline \texttt{unsym1} counts are zero in both AM and PM, and the tidal-candidate count is also zero.

The zero-count result is an empirical output of the current Shenzhen preprocessing and filter combination, not a structural statement about directional congestion in Shenzhen. The appropriate next diagnostic step is to inspect stage03b inclusion filters, effective directional pairing coverage by \texttt{cline\_id}, and peak-window support before interpreting this as evidence on the underlying asymmetry distribution.

\begin{table}[h!]
\centering
\caption{Stage03b asymmetry summary, Shenzhen run}
\label{tab:sz_stage03b_summary}
\begin{tabular}{lr}
\toprule
Statistic & Value \\
\midrule
Asymmetry rows & 0 \\
AM asymmetric centerlines (unsym1) & 0 \\
PM asymmetric centerlines (unsym1) & 0 \\
Reversed faster direction count & 0 \\
Tidal-lane candidates & 0 \\
\bottomrule
\end{tabular}
\end{table}

\begin{figure}[h!]
\centering
\includegraphics[width=0.72\textwidth]{../../data_work/outputs/shenzhen_test_line_v1_run1/figures/hist_speed_ratio_am_pm.png}
\caption{AM and PM speed-ratio histogram from the Shenzhen asymmetry diagnostic stage.}
\label{fig:sz_asym_ratio_hist}
\end{figure}

\section{Grid Construction}

Stage05 builds three grid systems on the Shenzhen boundary. The directed baseline backbone from stage01 is the network object used for within-cell diagnostics. This section now reports the smaller-grid Shenzhen run (\texttt{shenzhen\_smallgrid\_v1\_run1}).

\begin{table}[h!]
\centering
\caption{Stage05 grid comparison, Shenzhen run}
\label{tab:sz_stage05_grid}
\begin{tabular}{lrrrr}
\toprule
Grid type & Cells & Mean area (km$^2$) & Segment-within share & Length-within share \\
\midrule
Square & 695 & 3.30 & 65.55\% & 41.93\% \\
Hex & 1,698 & 0.83 & 25.56\% & 11.83\% \\
Voronoi & 505 & 4.54 & 59.67\% & 31.55\% \\
\bottomrule
\end{tabular}
\end{table}

\begin{figure}[h!]
\centering
\includegraphics[width=0.95\textwidth]{../../data_work/outputs/shenzhen_smallgrid_v1_run1/figures/aa_style_square3km_AM_2panel.png}
\caption{Square-grid AM diagnostics in the Shenzhen run: coverage and asymmetry panels.}
\label{fig:sz_square_am_2panel}
\end{figure}

\begin{figure}[h!]
\centering
\includegraphics[width=0.95\textwidth]{../../data_work/outputs/shenzhen_smallgrid_v1_run1/figures/aa_style_hex3km_AM_2panel.png}
\caption{Hex-grid AM diagnostics in the Shenzhen run: coverage and asymmetry panels.}
\label{fig:sz_hex_am_2panel}
\end{figure}

\begin{figure}[h!]
\centering
\includegraphics[width=0.95\textwidth]{../../data_work/outputs/shenzhen_smallgrid_v1_run1/figures/aa_style_voronoi3km_AM_2panel.png}
\caption{Voronoi-grid AM diagnostics in the Shenzhen run: coverage and asymmetry panels.}
\label{fig:sz_voronoi_am_2panel}
\end{figure}

\section{Grid-to-Grid Travel Costs and QSM Inputs}

\subsection{Stage06 travel-cost graph objects}

Stage06 outputs directional grid-link objects for all three grid systems. The Shenzhen run produces finite AM and PM edge tables in each partition.

\begin{table}[h!]
\centering
\caption{Stage06 network objects, Shenzhen run}
\label{tab:sz_stage06_summary}
\begin{tabular}{lrrrrr}
\toprule
Grid type & Links (long) & Links (agg) & Nodes & AM edges & PM edges \\
\midrule
Square & 3,322 & 1,408 & 328 & 704 & 704 \\
Hex & 4,013 & 2,213 & 437 & 1,106 & 1,107 \\
Voronoi & 4,043 & 2,099 & 304 & 1,049 & 1,050 \\
\bottomrule
\end{tabular}
\end{table}

\subsection{Stage07 OD reachability}

Stage07 maps commuters to grids and computes reachable OD pairs under the AM graph. Reachability is high in all three systems, with the strongest pair-level coverage in Voronoi.

\begin{table}[h!]
\centering
\caption{Stage07 OD reachability, Shenzhen run}
\label{tab:sz_stage07_summary}
\begin{tabular}{lrrrr}
\toprule
Grid type & OD pairs (all) & OD pairs (reachable) & Reachable pair share & Reachable pop share \\
\midrule
Square & 45,048 & 37,245 & 82.68\% & 98.42\% \\
Hex & 76,077 & 62,421 & 82.05\% & 96.97\% \\
Voronoi & 73,903 & 71,811 & 97.17\% & 99.65\% \\
\bottomrule
\end{tabular}
\end{table}

\subsection{Stage08 QSM exports}

Stage08 exports QSM-ready node, edge, and OD tables. In this Shenzhen draft, road commuters are defined as total commuters by configuration, so the QSM OD object does not rely on a mode-specific decomposition.

\begin{table}[h!]
\centering
\caption{Stage08 QSM export summary, Shenzhen run}
\label{tab:sz_stage08_summary}
\begin{tabular}{lrrrr}
\toprule
Grid type & QSM nodes & QSM edges & QSM OD pairs & Missing-lane edges \\
\midrule
Square & 305 & 704 & 37,245 & 0 \\
Hex & 413 & 1,106 & 62,421 & 0 \\
Voronoi & 304 & 1,049 & 71,811 & 0 \\
\bottomrule
\end{tabular}
\end{table}

\section{Implementation Notes for the Shenzhen Draft}

This draft is a Shenzhen counterpart to the Beijing manuscript structure in \texttt{TidalLanes\_0424.tex}, starting from the data-processing sections. The quantitative statements and figure paths are taken from \texttt{shenzhen\_test\_line\_v1\_run1} outputs.

[TODO: verify whether the Shenzhen asymmetry zero-count result is driven by filter design, directional pairing sparsity, or a preprocessing mismatch in stage03b.]

[TODO: lock the Shenzhen-specific figure list and update each section with figure references that match the final appendix layout.]

\end{document}
